\section{B. Inspection leads}\label{sec:b}

\begin{inshort}
Module B gives inspectors a short list of places to visit: new large buildings inside zones where no new
well permit can be issued, and outside the reach of piped water. A building like that must get water from
somewhere. Each visit is recorded in the app, which over time builds the per-site record that no open data
set has.
\end{inshort}

\subsection{Why not use sinking to find wells}

The first version of B ranked ``sinking that permits do not explain'' as a hunt for illegal wells. The
experiments made this unsound: sinking follows clay thickness and a regional pressure fall, and registered
wells sit on \emph{slower} ground \cite{exp13,exp14}. Unexplained sinking does not mean an illegal well. B
therefore uses building change and zoning. Sinking appears only as background, to show how damaged a zone
is.

\subsection{Facts that support the leads}

\begin{itemize}
\item \textbf{The zoning exists.} Dinas ESDM Jateng mapped the groundwater conservation zones of the CAT
  Semarang--Demak in 2022 under Permen ESDM 31/2018 \cite{undip,idx} \Off. The ``damaged'' (rusak) zone
  covers most of Kota Semarang. In critical and damaged zones no new abstraction permit can be issued;
  renewals get 50\,\% of the earlier volume and must build recharge wells.
\item \textbf{Industry depends on wells.} PDAM sold 0.22\Mm\ to industry in 2022 \cite{perdasmg} \Off.
\item \textbf{Unregistered wells are common.} In West Java, Dinas ESDM counted at least 6,000 well points,
  of which about 200 had a permit, and is remapping wells to raise tax \cite{ddtcjabar} \Wk. This is
  another province, from a news report.
\item \textbf{Satellites already find unrecorded abstraction elsewhere.} Jordan compared irrigated area from
  Landsat with recorded abstraction and found unlicensed wells and faulty meters \cite{jordan}. In Peru's
  Ica Valley, illicit use was about twice the regulated volume \cite{ica}. EU regulators rate optical and
  radar methods as mature but needing local tuning \cite{impel}.
\end{itemize}

\begin{caution}
All the foreign cases are irrigation, where crops show water use through evaporation. A factory well shows
nothing from space. B relies on a chain of reasoning (new, large, banned zone, no pipes, so probably a well),
and the chain must be tested (\V{4}, \V{5}).
\end{caution}

\subsection{Method \M{3}: the lead filter}

A building becomes a lead when it passes four filters. There is no weighted score: a score with made-up
weights would be a number with nothing behind it.

\begin{center}
\small
\begin{tabularx}{\linewidth}{@{}L{6mm}L{48mm}Y l@{}}
\toprule
 & \textbf{Filter} & \textbf{Data} & \textbf{Label}\\
\midrule
1 & Built after the 2022 zone map & GHSL and Landsat built-up change, building footprints \D{5} & \Mea\\
2 & Footprint above a size threshold, in a water-using class (factory, hotel, mall, hospital) & Footprints, OSM land use \D{5} & \Mea\\
3 & Inside a critical or damaged zone & Zone map \D{6} & \Off\\
4 & Outside the PDAM network & PDAM coverage per village (kelurahan) & \Mis\\
\bottomrule
\end{tabularx}
\end{center}

Filter 4 is the weakest link: the PDAM network map is not public. Until coverage is obtained (a PPID
information request to PDAM or the city), the app uses a proxy: distance to known mains and the SPAM
service districts. Leads based on the proxy are labelled as such.

\subsection{Method \M{4}: the field-feedback loop}

Each lead has a card with a satellite image, footprint, zone and PDAM status, and a button for the
inspector to record what was found:
\emph{registered well / unregistered well / no well / PDAM customer}. The records give two things:
\begin{itemize}
\item the hit rate of B per lead class, measured in the field (\V{6});
\item a growing per-site register, the data that files 12 and 13 showed is missing \cite{exp12,exp13}.
\end{itemize}

Lead cards are for inspectors only. The public view shows counts per district, not addresses: a lead is a
reason to visit, not an accusation (\R{7}).

\subsection{How B is checked}

\begin{itemize}
\item \V{4} \textbf{Arizona retro test.} Arizona registers every well with a drilling date and location
  (Wells 55), already used in file 05 \cite{exp05}. Do new large buildings outside water-provider service
  areas get a new registered well within 500\,m in the same or next two years, more often than
  built-matched random buildings? Pass marks are in Table~\ref{tab:valid}.
\item \V{5} \textbf{Semarang permit overlap.} Leads that fall on a known permitted well (Lo et al.\
  \cite{lo2022} or an ESDM permit list) are confirmed water users. Leads should hit permitted wells more
  often than random new buildings of the same size.
\item \V{6} \textbf{Field hit rate.} Reported per lead class with a 90\,\% interval. It has no pass mark:
  it is the number B exists to measure.
\end{itemize}

\begin{caution}
B cannot measure how many unpermitted wells it misses; that needs a full census. Factories that were already
self-supplied before 2022 are not ``new'' and are not leads. They belong to the declared volume in A.
\end{caution}
